% Materiał z AF 25; fragment do rozdziału o dystrybucjach temperowanych.
\subsection{Łańcuch przestrzeni sprzężonych i kryterium temperowania}

\begin{lemat}[Transpozycja gęstego zanurzenia]
\label{lem:transpozycja25}
Niech $X,Y$ będą przestrzeniami lokalnie wypukłymi
Hausdorffa, a $j:X\to Y$ ciągłym odwzorowaniem liniowym
o gęstym obrazie. Wówczas transpozycja
$j':Y'\to X'$, $j'(T)=T\circ j$, jest liniowa,
różnowartościowa i ciągła zarówno dla topologii słabych-*, jak i dla
silnych topologii przestrzeni dualnych.
\end{lemat}
\begin{proof}
Dla każdego $T\in Y'$ złożenie $T\circ j$ jest ciągłym
funkcjonałem liniowym na $X$, więc $j'$ jest dobrze
określone; liniowość wynika z liniowości składania
względem pierwszego czynnika. Jeżeli $T\circ j=0$,
to $T=0$ przez gęstość $j(X)$ i ciągłość $T$.
Dla $x\in X$ mamy
$|\langle j'T,x\rangle|=|\langle T,jx\rangle|$,
co daje ciągłość w topologiach słabych-*.
Jeśli $B\subseteq X$ jest ograniczony, to $j(B)$ jest
ograniczony w $Y$ (dla każdej ciągłej seminormy $p$
na $Y$ seminorma $p\circ j$ jest ciągła na $X$).
Dla seminorm silnych zachodzi dokładna równość
$\sup_{x\in B}|\langle j'T,x\rangle|
= \sup_{y\in j(B)}|\langle T,y\rangle|$.
To dowodzi ciągłości w topologiach silnych.
\end{proof}

\begin{wniosek}[Łańcuch przestrzeni dualnych]
\label{cor:dual-lancuch25}
Dla każdego $d\ge1$ zachodzą kanoniczne, ciągłe i
różnowartościowe odwzorowania liniowe
\[
 \mathcal E'(\R^d)\longrightarrow\mathcal S'(\R^d)
 \longrightarrow\mathcal D'(\R^d),
\]
zarówno w topologiach słabych-*, jak i silnych.
Dla każdego niepustego zbioru otwartego $\Omega\subseteq\R^d$
kanoniczne odwzorowanie
$\mathcal E'(\Omega)\to\mathcal D'(\Omega)$
ma te same własności. Samo stwierdzenie ciągłości i
różnowartościowości \emph{nie oznacza}, że są to zanurzenia
topologiczne na obrazy z topologią podprzestrzeni.
\end{wniosek}
\begin{proof}
Stosujemy Lemat~\ref{lem:transpozycja25} do odwzorowań
z Twierdzenia~\ref{thm:zanurzenia25}. Jeżeli $i:X\to Y$
oraz $j:Y\to Z$ są tymi kolejnymi odwzorowaniami,
to $(j\circ i)'=i'\circ j'$, gdyż dla $T\in Z'$
mamy $(T\circ j)\circ i=T\circ(j\circ i)$.
Dlatego ograniczenie z $\mathcal E'$ do $\mathcal D'$
przez $\mathcal S'$ jest tym samym odwzorowaniem co
bezpośrednie ograniczenie. W przypadku zespolonym
transpozycja jest liniowa, bez sprzężenia zespolonego.
Ciągłość i różnowartościowość nie zapewniają ciągłości
odwrotności na obrazie; poniżej podamy kontrprzykłady.
\end{proof}
\subsection{Dystrybucje temperowane jako podprzestrzeń \texorpdfstring{$\mathcal D'$}{D'}}

\begin{twierdzenie}[Kryterium temperowania]
\label{thm:temp-criterion25}
Niech $d\ge1$ będzie liczbą całkowitą, $\K\in\{\R,\C\}$
i $T\in\mathcal D'(\R^d)$ nad $\K$.
Następujące warunki są równoważne:
\begin{enumerate}[label=\textup{(\roman*)}]
\item $T$ należy do obrazu kanonicznego odwzorowania
$\mathcal S'(\R^d)\to\mathcal D'(\R^d)$;
\item istnieją $C>0$, $N,m\in\N_0$ takie, że
$|\langle T,\varphi\rangle|\le C p_{N,m}(\varphi)$
dla każdej $\varphi\in\mathcal D(\R^d)$.
\end{enumerate}
Jeżeli te warunki zachodzą, przedłużenie
$\widetilde T\in\mathcal S'(\R^d)$ jest jednoznaczne
i spełnia $|\langle\widetilde T,f\rangle|\le C p_{N,m}(f)$
dla wszystkich $f\in\mathcal S(\R^d)$, z tymi samymi
stałymi co w~\textup{(ii)}. W szczególności każda
dystrybucja o zwartym nośniku jest temperowana.
\end{twierdzenie}
\begin{proof}
Niech najpierw $\widetilde T\in\mathcal S'$ będzie
przedłużeniem $T$. Z ciągłości w zerze i współkońcowości
rosnącego ciągu norm $p_{j,j}$ istnieją $j\in\N_0$ oraz
$\delta>0$ takie, że $p_{j,j}(f)<\delta$ implikuje
$|\widetilde T(f)|<1$. Dla $f\ne0$ stosujemy ten warunek
do $\delta f/(2p_{j,j}(f))$, otrzymując
$|\widetilde T(f)|\le (2/\delta)p_{j,j}(f)$.
Dla $f=0$ nierówność jest oczywista. Ograniczenie
do $\mathcal D$ daje warunek~\textup{(ii)}.
Jest to argument z Twierdzenia~\ref{thm:tempered-estimate22},
działający nad obydwoma ciałami.

Załóżmy teraz~\textup{(ii)}. Dla każdego $f\in\mathcal S$
wybieramy ciąg $(\varphi_j)$ w $\mathcal D$ zbieżny do $f$
w $\mathcal S$ (Twierdzenie~\ref{thm:zanurzenia25}).
Ponieważ
\[
 |T(\varphi_j)-T(\varphi_k)|
 \le C p_{N,m}(\varphi_j-\varphi_k)\longrightarrow0
 \quad(j,k\to\infty),
\]
ciąg wartości ma granicę w zupełnym ciele $\K$.
Kładziemy $\widetilde T(f)=\lim_jT(\varphi_j)$.
Dla drugiego ciągu $(\psi_j)$ zbieżnego do $f$ mamy
$|T(\varphi_j)-T(\psi_j)|\le
C p_{N,m}(\varphi_j-\psi_j)\to0$, więc definicja nie zależy
od wyboru ciągu. Dla $f,g\in\mathcal S$ i $a,b\in\K$
można aproksymować $af+bg$ przez $a\varphi_j+b\psi_j$,
gdzie teraz $\varphi_j\to f$ i $\psi_j\to g$.
Przechodząc do granicy, otrzymujemy liniowość.
Ponadto $|p_{N,m}(\varphi_j)-p_{N,m}(f)|
\le p_{N,m}(\varphi_j-f)$, więc
\[
 |\widetilde T(f)|\le C p_{N,m}(f).
\]
Daje to ciągłość; na $\mathcal D$ wystarczy wziąć
ciąg stały, aby uzyskać $\widetilde T|_{\mathcal D}=T$.
Jednoznaczność wynika z gęstości $\mathcal D$.

Jeżeli $T$ ma zwarty nośnik, to z
Twierdzenia~\ref{thm:compact-support-E25} otrzymujemy
$|T(\varphi)|\le C p_{K,m}(\varphi)\le C p_{0,m}(\varphi)$
na funkcjach testowych, a zatem~\textup{(ii)} zachodzi z $N=0$.
Przedłużenie jest ograniczeniem funkcjonału na $\mathcal E$
do $\mathcal S$, przez jednoznaczność; odpowiada to
Stwierdzeniu~\ref{prop:schwartz-embeddings22}.
\end{proof}

\begin{przyklad}[Dystrybucja nietemperowana]
\label{ex:non-tempered25}
Na $\R$ funkcja $f(x)=e^x$ jest lokalnie całkowalna,
więc wyznacza dystrybucję regularną $T_f\in\mathcal D'(\R)$.
Nie jest ona temperowana. Wybierzmy nieujemną, niezerową
$\psi\in\mathcal D((-1,1))$, przedłużoną zerem na $\R$,
i niech $\varphi_n(x)=e^{-n/2}\psi(x-n)$ dla $n\ge1$.
Dla $x\in\supp\varphi_n$ mamy $|x|\le n+1$; pochodne
$\varphi_n$ znikają poza tym nośnikiem. Stąd dla wszystkich
$N,m\in\N_0$
\[
 p_{N,m}(\varphi_n)\le e^{-n/2}(n+2)^N p_{0,m}(\psi)
 \longrightarrow0,
\]
gdyż funkcja wykładnicza rośnie szybciej niż każdy
wielomian. Jednak po podstawieniu $y=x-n$ otrzymujemy
\[
 \langle T_f,\varphi_n\rangle=
 e^{n/2}\int_\R e^y\psi(y)\,\mathrm d\lambda_1(y)
 \longrightarrow+\infty.
\]
Całka jest skończona i ściśle dodatnia: $\psi$ jest
nieujemna i dodatnia na pewnym niepustym przedziale.
Ciągłość przedłużenia na $\mathcal S$ wymagałaby
$\langle T_f,\varphi_n\rangle\to0$, co daje sprzeczność.
Przykład działa nad $\R$ i $\C$; razem z
Przykładem~\ref{ex:non-tempered22} pokazuje on dwie
metody wykrywania braku temperowalności (przesuwanie
nośnika oraz powiększanie odcięcia).
\end{przyklad}
\begin{przyklad}[Zbieżność i ucieczka nośnika]
\label{ex:escaping-deltas25}
Na $\R$ niech $\delta_n$ oznacza deltę Diraca w punkcie
$n$ dla $n\ge1$. Każda z tych dystrybucji ma zwarty
nośnik, więc należy także do $\mathcal E'$ i $\mathcal S'$
na mocy Twierdzeń~\ref{thm:compact-support-E25}
i~\ref{thm:temp-criterion25}.
Wtedy $\delta_n\to0$ silnie w $\mathcal D'(\R)$:
dla ograniczonego $B\subseteq\mathcal D(\R)$ nośniki
wszystkich jego elementów mieszczą się w jednym zwartym
przedziale (Twierdzenie~\ref{thm:D-bounded18}), więc
$\sup_{\varphi\in B}|\varphi(n)|=0$ dla dużych $n$.
Również $\delta_n\to0$ silnie w $\mathcal S'(\R)$, bo dla
ograniczonego zbioru $B\subseteq\mathcal S$ mamy
$C_B=\sup_{\varphi\in B}p_{2,0}(\varphi)<\infty$ oraz
$\sup_{\varphi\in B}|\varphi(n)|\le C_B(1+n)^{-2}$.
Jednak $\delta_n\not\to0$ nawet słabo-* w
$\mathcal E'(\R)$, ponieważ $\langle\delta_n,1\rangle=1$.

Z kolei $T_n=e^n\delta_n$ zbiega silnie do zera w
$\mathcal D'(\R)$ (na każdym ograniczonym zbiorze
seminorma silna jest od pewnego indeksu zerowa).
Nie zbiega do zera nawet słabo-* w $\mathcal S'(\R)$.
Niech $r(x)=\sqrt{1+x^2}$ i $\varphi(x)=e^{-r(x)/2}$.
Wykażemy, że $\varphi\in\mathcal S(\R)$.
Dla każdego $k\in\N_0$ istnieje wielomian $P_k(x,z)$
taki, że
\[
 \varphi^{(k)}(x)=e^{-r(x)/2}r(x)^{-2k}P_k(x,r(x)).
\]
Dla $k=0$ przyjmujemy $P_0=1$. Ponieważ $r'=x/r$,
różniczkowanie wzoru daje rekurencję
\[
 P_{k+1}=z^2\partial_x P_k+xz\partial_z P_k
          -2kxP_k-\tfrac12 xzP_k.
\]
Zatem wszystkie $P_k$ są wielomianami. Ponieważ
$1\le r(x)\le1+|x|$ i $r(x)\ge|x|$, każda pochodna
jest ograniczona przez stałą razy wielomianowy wzrost
$(1+|x|)^{M_k}$ pomnożony przez $e^{-|x|/2}$.
Wszystkie seminormy Schwartza są więc skończone.
Jednak dla $n\ge1$ mamy $r(n)\le n+1$, toteż
$\langle T_n,\varphi\rangle=e^{n-r(n)/2}
\ge e^{(n-1)/2}\to+\infty$.

Dla dowolnego niepustego zbioru otwartego $\Omega$
wybierzmy wyczerpanie $(K_j)$ i punkty
$a_j\in\Omega\setminus K_j$.
Wówczas $\delta_{a_j}\to0$ silnie w $\mathcal D'(\Omega)$,
gdyż każdy ograniczony zbiór testów ma wspólny zwarty
nośnik zawarty w pewnym $K_J$.
Natomiast $\langle\delta_{a_j},1\rangle=1$ na
$\mathcal E(\Omega)$ dla wszystkich $j$.
Zatem wszystkie odwzorowania z
Wniosku~\ref{cor:dual-lancuch25} nie są izomorfizmami
topologicznymi na swoje obrazy z topologią podprzestrzeni,
ani dla topologii słabych-*, ani dla silnych.
\end{przyklad}

\subsection{Zgodność operacji i granice ich działania}

\begin{stwierdzenie}[Różniczkowanie i mnożenie]
\label{prop:ops25}
Niech $\Omega\subseteq\R^d$ będzie niepustym zbiorem otwartym.
Dla każdego wielowskaźnika $\alpha\in\N_0^d$
operator $\partial^\alpha$ działa liniowo i ciągle na
$\mathcal E(\Omega)$ i $\mathcal D(\Omega)$, a przez
transpozycję także na $\mathcal E'(\Omega)$ i
$\mathcal D'(\Omega)$ w topologiach słabych-* i silnych.
Dla $a\in C^\infty(\Omega)$ mnożenie $f\mapsto af$
i odpowiadające mu mnożenie dystrybucji działają ciągle
na tych przestrzeniach.
Różniczkowanie działa liniowo i ciągle również na
$\mathcal S(\R^d)$ i na $\mathcal S'(\R^d)$ w obu
topologiach dualnych. Na tych przestrzeniach
te same własności mnożenia zachodzą dla
$a\in C^\infty(\R^d)$, jeżeli dla każdego
wielowskaźnika $\alpha$ istnieją $C_\alpha>0$ i
$N_\alpha\in\N_0$ takie, że
$|\partial^\alpha a(x)|\le C_\alpha(1+|x|)^{N_\alpha}$
dla wszystkich $x\in\R^d$.
\end{stwierdzenie}
\begin{proof}
Dla $\mathcal E$ i $\mathcal D$ oszacowania seminorm
dla różniczkowania i mnożenia wynikają z reguły Leibniza,
jak w stwierdzeniu~\ref{prop:operations-cinf17}
i twierdzeniach~\ref{thm:dist-deriv-cont20}, \ref{thm:multiply20}.
Dla $\mathcal S$ i $\mathcal S'$ stosujemy stwierdzenie
\ref{prop:schwartz-operators17} oraz twierdzenie
\ref{thm:tempered-operators22}, w tym oszacowanie
mnożenia przez funkcję o pochodnych wzrostu wielomianowego.
Dla ciągłego operatora $A:X\to X$ transpozycja spełnia
$P_B(A'T)=P_{A(B)}(T)$, gdzie
$P_B(T)=\sup_{f\in B}|T(f)|$ i $A(B)$ jest ograniczony.
Daje to ciągłość silną, a dla singletonów także słabą-*.
Stosujemy ten argument w szczególności do
$X=\mathcal E$ i operatorów $(-1)^{|\alpha|}\partial^\alpha$
oraz $f\mapsto af$. Operatory na funkcjach komutują
z kanonicznymi odwzorowaniami przestrzeni funkcji,
więc przez transpozycję odpowiadające operacje
na przestrzeniach sprzężonych komutują z ograniczeniami
z wniosku~\ref{cor:dual-lancuch25} (w przypadku $\mathcal S'$
przy podanym warunku wzrostu mnożnika).
\end{proof}

\begin{przyklad}[Nie każdy mnożnik gładki zachowuje przestrzeń Schwartza]
\label{ex:not-multiplier25}
Funkcja $a(x)=e^{|x|^2}$ jest gładka, ale mnożenie przez nią
nie zachowuje przestrzeni Schwartza: dla
$f(x)=e^{-|x|^2}\in\mathcal S(\R^d)$ iloczyn $af=1$
nie należy do $\mathcal S(\R^d)$. Przynależność funkcji Gaussa do $\mathcal S$ wynika
z tego, że każda jej pochodna jest wielomianem
pomnożonym przez tę funkcję (indukcja z reguły
iloczynu), a spadek Gaussa dominuje każdą potęgę.
Funkcja $1$ nie spełnia warunku $p_{1,0}(1)<\infty$.
W przestrzeniach $\mathcal D$ i $\mathcal E$
mnożenie przez $a$ jest natomiast ciągłe.
Co więcej, zwykła dystrybucja $aT_1=T_a$ nie jest
temperowana, choć $T_1$ jest temperowana
(Przykłady~\ref{ex:tempered-regular22}
i~\ref{ex:non-tempered22}).
\end{przyklad}

\begin{uwaga}[Transformacja Fouriera i splot]
\label{rem:fourier-conv25}
Nad $\C$ transformacja Fouriera jest automorfizmem
liniowo-topologicznym $\mathcal S(\R^d)$ i
$\mathcal S'(\R^d)$ w obu topologiach dualnych
(Twierdzenia~\ref{thm:fourier-schwartz17}
i~\ref{thm:fourier-auto-tempered22}).
Nie jest zdefiniowana przez tę samą transpozycję
na całej $\mathcal D'$: niezerowa funkcja testowa
nie może mieć transformaty Fouriera o zwartym nośniku.
Istotnie, dla $f\in\mathcal D$ jej transformata
jest rzeczywiście analityczna: rozwinięcie wykładnika
w szereg potęgowy w otoczeniu dowolnego punktu
można całkować wyraz po wyrazie, ponieważ $f$ ma
zwarty nośnik, a szereg jest jednostajnie bezwzględnie
zbieżny na nim dla parametrów z każdego ograniczonego
otoczenia. Gdyby transformata miała zwarty nośnik,
znikałaby na niepustym otwartym podzbiorze $\R^d$,
a z jednoznaczności funkcji rzeczywiście analitycznej
znikałaby na całym spójnym $\R^d$.
Różnowartościowość transformaty na $\mathcal S$
dawałaby $f=0$. Zatem wyrażenie
$\langle T,\mathcal F f\rangle$ nie musi mieć sensu
dla $T\in\mathcal D'$ i $f\in\mathcal D$.

Dla ustalonego $T\in\mathcal E'(\R^d)$ splot
$S\mapsto T*S$ jest ciągłym operatorem liniowym
na $\mathcal D'(\R^d)$ w obu topologiach
(Twierdzenie~\ref{thm:convolution-properties21}).
Aby sprawdzić, że zachowuje również $\mathcal S'$,
niech $L$ będzie zwartym zbiorem i $C>0$, $m\in\N_0$
stałymi oszacowania przedłużenia $T$ na $\mathcal E$
z Twierdzenia~\ref{thm:compact-support-E25}.
Dla $f\in\mathcal S$ określamy
$(A_Tf)(y)=\langle T_x,f(x+y)\rangle$,
korzystając z przedłużenia na $\mathcal E$.
Z ciągłości $T$ na $\mathcal E$ i różniczkowania
przesunięć w jej topologii otrzymujemy
$\partial_y^\beta A_T f(y)=
\langle T_x,\partial^\beta f(x+y)\rangle$:
ilorazy różnicowe wraz z pochodnymi względem $x$
zbiegają jednostajnie na każdym zwartym zbiorze
przez wzór całkowy na różnicę funkcji gładkiej.
Ponieważ $1+|y|\le(1+|x|)(1+|x+y|)$, dla
$R=\max_{x\in L}|x|$ otrzymujemy
\[
 p_{N,k}(A_Tf)\le C(1+R)^N p_{N,k+m}(f).
\]
Zatem $A_T:\mathcal S\to\mathcal S$ jest liniowe
i ciągłe. Jego transpozycja
$\langle T*S,f\rangle=\langle S,A_Tf\rangle$
jest ciągła na $\mathcal S'$ w obu topologiach
i zgadza się na $\mathcal D$ z definicją splotu
z rozdziału~\ref{sec:splot21}.
Nie wynika stąd istnienie splotu dwóch dowolnych
dystrybucji (Uwaga~\ref{rem:convolution-limits21}).
\end{uwaga}
